% LTeX: enabled=false
% Sección 3: correspondencia entre objetivos y resultados, y próximos pasos

\seccion{Balance}

% ===============================================================
% 12. Objetivos y resultados
% ===============================================================

\begin{frame}{Lo que nos propusimos y lo que logramos}
    \begin{tabla}{colspec = {Q[c,m,0.7cm] Q[l,m,4.3cm] X[l,m] Q[c,m,2.3cm]},
                  rows = {font=\footnotesize}, rowsep = 1pt}
        Obj. & Qué nos propusimos & Qué logramos & Estado \\
        1 & Analizar un conjunto de datos público & CornOrb seleccionado y revisado, con los datos a corregir identificados & \etiq{verde}{Alcanzado} \\
        2 & Construir el modelo de IA & Primera prueba con los mapas: detecta el 92\,\% de los queratoconos sin falsos positivos. Se siguen probando redes y estrategias, y falta la fusión con los datos clínicos & \etiq{azul}{En proceso} \\
        3 & Integrar la explicación visual & Grad-CAM funcionando sobre los mapas. Falta mostrarlo en el informe & \etiq{azul}{En proceso} \\
        4 & Desarrollar la aplicación web & Partes y herramientas definidas (React y FastAPI) & \etiq{azul}{En proceso} \\
        --- & Conseguir datos locales para validar & 9 centros contactados, charla con Roberto Albertazzi y trámites en curso & \etiq{gray}{Por empezar} \\
    \end{tabla}
\end{frame}

% ===============================================================
% 13. Próximos pasos
% ===============================================================
\begin{frame}{Próximos pasos}
    % Noviembre y Diciembre lado a lado, con cuatro puntos cada uno para que
    % los bloques queden de la misma altura. "En paralelo" abajo, a todo el ancho
    % de los dos. Las columnas ocupan 0,99 (Beamer reparte el espacio sobrante
    % entre ellas), y el bloque de abajo mide lo mismo para que los bordes coincidan
    \vspace{-0.5cm}
    \begin{columns}[T, totalwidth=0.99\textwidth]
        \begin{column}{0.47\textwidth}
            \begin{block}{Noviembre}
                \begin{itemize}
                    \item Seguir probando redes y estrategias.
                    \item Combinar los mapas con los datos clínicos en un solo modelo.
                    \item Ajustarlo para detectar al menos el 95\,\% de los queratoconos.
                    \item Empezar la aplicación web.
                \end{itemize}
            \end{block}
        \end{column}
        \hspace{0.03\textwidth}
        \begin{column}{0.47\textwidth}
            \begin{block}{Diciembre}
                \begin{itemize}
                    \item Terminar la aplicación web.
                    \item Conectar el modelo con la aplicación.
                    \item Generar el informe con el resultado y el mapa de calor.
                    \item Tener el prototipo funcionando de punta a punta.
                \end{itemize}
            \end{block}
        \end{column}
    \end{columns}
    \vspace{0.3cm}
    % Sin center: las columnas arrancan en el margen izquierdo, así que este
    % bloque también, para que los bordes coincidan
    \noindent\begin{minipage}{0.99\textwidth}
        \begin{block}{En paralelo}
            Avanzar con los trámites para conseguir los datos locales de validación.
        \end{block}
    \end{minipage}
\end{frame}
